\documentclass[a4paper,12pt,twoside]{memoir}

% Castellano
\usepackage[spanish,es-tabla]{babel}
\selectlanguage{spanish}
\usepackage[utf8]{inputenc}
\usepackage[T1]{fontenc}
\usepackage{lmodern} % Scalable font
\usepackage{microtype}
\usepackage{placeins}

\RequirePackage{booktabs}
\RequirePackage[table]{xcolor}
\RequirePackage{xtab}
\RequirePackage{multirow}

% Links
\PassOptionsToPackage{hyphens}{url}\usepackage[colorlinks]{hyperref}
\hypersetup{
	allcolors = {red}
}

% Ecuaciones
\usepackage{amsmath}

% Rutas de fichero / paquete
\newcommand{\ruta}[1]{{\sffamily #1}}

% Párrafos
\nonzeroparskip

% Huérfanas y viudas
\widowpenalty100000
\clubpenalty100000

% Imagenes
\usepackage{graphicx}
\newcommand{\imagen}[2]{
	\begin{figure}[!h]
		\centering
		\includegraphics[width=0.9\textwidth]{#1}
		\caption{#2}\label{fig:#1}
	\end{figure}
	\FloatBarrier
}

\newcommand{\imagenflotante}[2]{
	\begin{figure}%[!h]
		\centering
		\includegraphics[width=0.9\textwidth]{#1}
		\caption{#2}\label{fig:#1}
	\end{figure}
}

\newcommand{\imagenMediana}[2]{
	\begin{figure}[!h]
		\centering
		\includegraphics[width=0.6\textwidth]{#1}
		\caption{#2}\label{fig:#1}
	\end{figure}
	\FloatBarrier
}

\newcommand{\imagenPequena}[2]{
	\begin{figure}[!h]
		\centering
		\includegraphics[width=0.2\textwidth]{#1}
		\caption{#2}\label{fig:#1}
	\end{figure}
	\FloatBarrier
}



% El comando \figura nos permite insertar figuras comodamente, y utilizando
% siempre el mismo formato. Los parametros son:
% 1 -> Porcentaje del ancho de página que ocupará la figura (de 0 a 1)
% 2 --> Fichero de la imagen
% 3 --> Texto a pie de imagen
% 4 --> Etiqueta (label) para referencias
% 5 --> Opciones que queramos pasarle al \includegraphics
% 6 --> Opciones de posicionamiento a pasarle a \begin{figure}
\newcommand{\figuraConPosicion}[6]{%
  \setlength{\anchoFloat}{#1\textwidth}%
  \addtolength{\anchoFloat}{-4\fboxsep}%
  \setlength{\anchoFigura}{\anchoFloat}%
  \begin{figure}[#6]
    \begin{center}%
      \Ovalbox{%
        \begin{minipage}{\anchoFloat}%
          \begin{center}%
            \includegraphics[width=\anchoFigura,#5]{#2}%
            \caption{#3}%
            \label{#4}%
          \end{center}%
        \end{minipage}
      }%
    \end{center}%
  \end{figure}%
}

%
% Comando para incluir imágenes en formato apaisado (sin marco).
\newcommand{\figuraApaisadaSinMarco}[5]{%
  \begin{figure}%
    \begin{center}%
    \includegraphics[angle=90,height=#1\textheight,#5]{#2}%
    \caption{#3}%
    \label{#4}%
    \end{center}%
  \end{figure}%
}
% Para las tablas
\newcommand{\otoprule}{\midrule [\heavyrulewidth]}
%
% Nuevo comando para tablas pequeñas (menos de una página).
\newcommand{\tablaSmall}[5]{%
 \begin{table}
  \begin{center}
   \rowcolors {2}{gray!35}{}
   \begin{tabular}{#2}
    \toprule
    #4
    \otoprule
    #5
    \bottomrule
   \end{tabular}
   \caption{#1}
   \label{tabla:#3}
  \end{center}
 \end{table}
}

%
% Nuevo comando para tablas pequeñas (menos de una página).
\newcommand{\tablaSmallSinColores}[5]{%
 \begin{table}[H]
  \begin{center}
   \begin{tabular}{#2}
    \toprule
    #4
    \otoprule
    #5
    \bottomrule
   \end{tabular}
   \caption{#1}
   \label{tabla:#3}
  \end{center}
 \end{table}
}

\newcommand{\tablaApaisadaSmall}[5]{%
\begin{landscape}
  \begin{table}
   \begin{center}
    \rowcolors {2}{gray!35}{}
    \begin{tabular}{#2}
     \toprule
     #4
     \otoprule
     #5
     \bottomrule
    \end{tabular}
    \caption{#1}
    \label{tabla:#3}
   \end{center}
  \end{table}
\end{landscape}
}

%
% Nuevo comando para tablas grandes con cabecera y filas alternas coloreadas en gris.
\newcommand{\tabla}[6]{%
  \begin{center}
    \tablefirsthead{
      \toprule
      #5
      \otoprule
    }
    \tablehead{
      \multicolumn{#3}{l}{\small\sl continúa desde la página anterior}\\
      \toprule
      #5
      \otoprule
    }
    \tabletail{
      \hline
      \multicolumn{#3}{r}{\small\sl continúa en la página siguiente}\\
    }
    \tablelasttail{
      \hline
    }
    \bottomcaption{#1}
    \rowcolors {2}{gray!35}{}
    \begin{xtabular}{#2}
      #6
      \bottomrule
    \end{xtabular}
    \label{tabla:#4}
  \end{center}
}

%
% Nuevo comando para tablas grandes con cabecera.
\newcommand{\tablaSinColores}[6]{%
  \begin{center}
    \tablefirsthead{
      \toprule
      #5
      \otoprule
    }
    \tablehead{
      \multicolumn{#3}{l}{\small\sl continúa desde la página anterior}\\
      \toprule
      #5
      \otoprule
    }
    \tabletail{
      \hline
      \multicolumn{#3}{r}{\small\sl continúa en la página siguiente}\\
    }
    \tablelasttail{
      \hline
    }
    \bottomcaption{#1}
    \begin{xtabular}{#2}
      #6
      \bottomrule
    \end{xtabular}
    \label{tabla:#4}
  \end{center}
}

%
% Nuevo comando para tablas grandes sin cabecera.
\newcommand{\tablaSinCabecera}[5]{%
  \begin{center}
    \tablefirsthead{
      \toprule
    }
    \tablehead{
      \multicolumn{#3}{l}{\small\sl continúa desde la página anterior}\\
      \hline
    }
    \tabletail{
      \hline
      \multicolumn{#3}{r}{\small\sl continúa en la página siguiente}\\
    }
    \tablelasttail{
      \hline
    }
    \bottomcaption{#1}
  \begin{xtabular}{#2}
    #5
   \bottomrule
  \end{xtabular}
  \label{tabla:#4}
  \end{center}
}



\definecolor{cgoLight}{HTML}{EEEEEE}
\definecolor{cgoExtralight}{HTML}{FFFFFF}

%
% Nuevo comando para tablas grandes sin cabecera.
\newcommand{\tablaSinCabeceraConBandas}[5]{%
  \begin{center}
    \tablefirsthead{
      \toprule
    }
    \tablehead{
      \multicolumn{#3}{l}{\small\sl continúa desde la página anterior}\\
      \hline
    }
    \tabletail{
      \hline
      \multicolumn{#3}{r}{\small\sl continúa en la página siguiente}\\
    }
    \tablelasttail{
      \hline
    }
    \bottomcaption{#1}
    \rowcolors[]{1}{cgoExtralight}{cgoLight}

  \begin{xtabular}{#2}
    #5
   \bottomrule
  \end{xtabular}
  \label{tabla:#4}
  \end{center}
}


















\graphicspath{ {./img/} }

% Capítulos
\chapterstyle{bianchi}
\newcommand{\capitulo}[2]{
	\setcounter{chapter}{#1}
	\setcounter{section}{0}
	\setcounter{figure}{0}
	\setcounter{table}{0}
	\chapter*{#2}
	\addcontentsline{toc}{chapter}{#2}
	\markboth{#2}{#2}
}

% Apéndices
\renewcommand{\appendixname}{Apéndice}
\renewcommand*\cftappendixname{\appendixname}

\newcommand{\apendice}[1]{
	%\renewcommand{\thechapter}{A}
	\chapter{#1}
}

\renewcommand*\cftappendixname{\appendixname\ }

% Formato de portada
\makeatletter
\usepackage{xcolor}
\newcommand{\tutor}[1]{\def\@tutor{#1}}
\newcommand{\course}[1]{\def\@course{#1}}
\definecolor{cpardoBox}{HTML}{E6E6FF}
\def\maketitle{
  \null
  \thispagestyle{empty}
  % Cabecera ----------------
\noindent\includegraphics[width=\textwidth]{cabecera}\vspace{1cm}%
  \vfill
  % Título proyecto y escudo informática ----------------
  \colorbox{cpardoBox}{%
    \begin{minipage}{.8\textwidth}
      \vspace{.5cm}\Large
      \begin{center}
      \textbf{TFG del Grado en Ingeniería Informática}\vspace{.6cm}\\
      \textbf{\LARGE\@title{}}\vspace{.6cm}\\
      \textbf{UBU-Chatbot}
      \end{center}
      \vspace{.2cm}
    \end{minipage}

  }%
  \hfill\begin{minipage}{.20\textwidth}
    \includegraphics[width=\textwidth]{escudoInfor}
  \end{minipage}
  \vfill
  % Datos de alumno, curso y tutores ------------------
  \begin{center}%
  {%
    \noindent\LARGE
    Presentado por \@author{}\\ 
    en Universidad de Burgos --- julio de 2021\\
    Tutor: \@tutor{}\\
  }%
  \end{center}%
  \null
  \cleardoublepage
  }
\makeatother

\newcommand{\nombre}{Nombre del alumno} %%% cambio de comando

% Datos de portada
\title{Desarrollo y explotación de un chatbot de FAQs sobre una asignatura de Trabajo Fin de Grado}
\author{Alfredo Asensio Vázquez}
\tutor{Dr. Raúl Marticorena Sánchez}
\date{\today}

\begin{document}

\maketitle

\renewcommand*\abstractname{Agradecimientos}
\begin{abstract}
	Este trabajo supone la culminación de muchos años de trabajo y esfuerzo. No ha sido un camino fácil ni ha estado exento de dificultades, pero nunca he caminado solo, ya que han sido muchas las personas que me han acompañado y a las que estoy muy agradecido y espero haber sido capaz de transmitírselo así.
	
	Pero quiero dedicar este logro que supone acabar Ingeniería a quienes más han sufrido con ello y siempre han estado ahí: mis padres.
	Gracias por todo vuestro sacrificio, apoyo y confianza. En ningún momento dudasteis de que lo conseguiría y el tiempo os ha dado la razón. Sin vosotros esto no hubiera sido posible y nada me hace más feliz que veros orgullosos de mí.
	
	Sois los mejores padres que podría tener,
	
		
	Os quiero.
\end{abstract}

 


\newpage\null\thispagestyle{empty}\newpage



%%%%%%%%%%%%%%%%%%%%%%%%%%%%%%%%%%%%%%%%%%%%%%%%%%%%%%%%%%%%%%%%%%%%%%%%%%%%%%%%%%%%%%%%
\thispagestyle{empty}


\noindent\includegraphics[width=\textwidth]{cabecera}\vspace{1cm}

\noindent D. Raúl Marticorena Sánchez, profesor del departamento de Ingeniería Informática, área de Lenguajes y Sistemas Informáticos.

\noindent Expone:

\noindent Que el alumno D. Alfredo Asensio Vázquez, con DNI 73007684A, ha realizado el Trabajo final de Grado en Ingeniería Informática titulado "Desarrollo y explotación de un chatbot de FAQs sobre una asignatura de Trabajo Fin de Grado
". 

\noindent Y que dicho trabajo ha sido realizado por el alumno bajo la dirección del que suscribe, en virtud de lo cual se autoriza su presentación y defensa.

\begin{center} %\large
En Burgos, {\large \today}
\end{center}

\vfill\vfill\vfill

% Author and supervisor
\begin{minipage}{0.45\textwidth}
\begin{flushleft} %\large
Vº. Bº. del Tutor:\\[2cm]
D. Raúl Marticorena Sánchez
\end{flushleft}
\end{minipage}
\hfill
%\begin{minipage}{0.45\textwidth}
%\begin{flushleft} %\large
%Vº. Bº. del co-tutor:\\[2cm]
%D. nombre co-tutor
%\end{flushleft}
%\end{minipage}
%\hfill

\vfill

% para casos con solo un tutor comentar lo anterior
% y descomentar lo siguiente
%Vº. Bº. del Tutor:\\[2cm]
%D. nombre tutor


\newpage\null\thispagestyle{empty}\newpage



\frontmatter

% Abstract en castellano
\renewcommand*\abstractname{Resumen}
\begin{abstract}
Los estudiantes que realizan el TFG a menudo tienen dudas acerca de su organización y funcionamiento. Muchas de estas dudas están ya recogidas en la documentación de la plataforma Moodle UBUVirtual o en los distintos documentos relacionados del portal web de la Universidad.

En este trabajo, se desarrolla un bot conversacional (chatbot) cuyo objetivo es dar respuesta a todas estas dudas que los alumnos puedan tener sobre el funcionamiento de la asignatura  Trabajo de Fin de Grado, reduciendo así también la carga de trabajo de los docentes.

Es también una de las primeras aproximaciones de la UBU a la implementación de una herramienta en auge como son los chatbots, que pueden aportar mucho valor en una futura explotación y expansión en las distintas plataformas que utiliza la Universidad. 

El ámbito del trabajo se centra principalmente en la integración del chatbot en la plataforma Moodle de la Universidad, llamada UBUVirtual, aunque no se limita a esta única plataforma, ya que también se realizará una integración con la plataforma Slack. 


\end{abstract}

\renewcommand*\abstractname{Descriptores}
\begin{abstract}
Chatbot, bot conversacional, agente inteligente, Dialogflow, Trabajo Fin de Grado, Moodle, Slack.
\end{abstract}

\clearpage

% Abstract en inglés
\renewcommand*\abstractname{Abstract}
\begin{abstract}
Students who take the End-of-Degree project often have doubts about its organization and operation. Many of these doubts are already collected in the documentation of the Moodle UBUVirtual platform or in the different related documents of the University web portal.

In this work we develop a chatbot whose objective is to answer all these doubts that students may have about the End of Degree Project subject, thus also reducing the workload of teachers.

It is also one of the UBU's first approaches to the implementation of a booming tool such as chatbots, which can add a lot of value in future exploitation and expansion on the different platforms used by the University. . 

The scope of the work is mainly focused on the integration of the chatbot in the University's Moodle platform, called UBUVirtual, although it is not limited to this single platform, since an integration with the Slack platform will also be carried out. 
\end{abstract}

\renewcommand*\abstractname{Keywords}
\begin{abstract}
Chatbot, intelligent agent, Dialogflow, End-of-Degree Project, Moodle, Slack.
\end{abstract}

\clearpage

% Indices
\tableofcontents

\clearpage

\listoffigures

\clearpage

\listoftables
\clearpage

\mainmatter
\capitulo{1}{Introducción}

\section{Contenido del trabajo}

Llegado al último curso de grado todos los alumnos de la UBU deben realizar su Trabajo de Fin de Grado (TFG) y son muchas las dudas que tienen acerca del funcionamiento del mismo, especialmente en la enseñanza online en la que al no haber un contacto tan estrecho con profesores y otros alumnos se tiene menos información. Motivo por el que los alumnos, aparte de consultar la extensa documentación aportada en las distintas plataformas de la UBU consultan a los profesores por email para despejar sus dudas.

Muchas de estas preguntas acerca del trabajo ya han sido respondidas anteriormente, o están explicadas en alguna parte de la documentación de UBUVirtual o de la página del grado. Con el propósito de ayudar a los alumnos a despejar estas dudas y reducir la carga de preguntas que reciben los profesores, se desarrolla este chatbot cuya finalidad es la de dar respuesta de forma automática a las preguntas de los usuarios acerca de todos los aspectos relacionados con la asignatura. Como restricción ha de utilizar únicamente herramientas gratuitas. El chatbot está integrado en la plataforma de UBUVirtual de la asignatura de TFG, a la que todos los alumnos matriculados en el trabajo tienen acceso. 

Por medio del aprendizaje automático nuestro chatbot será capaz de entender y responder las preguntas que los alumnos le formulen, y a su vez aprenderá de aquellas que no las haya recibido anteriormente para ser capaz de responderlas en la próxima ocasión en que se le pregunten.\\
Con ello resolveremos la doble problemática para profesores y alumnos, con la ventaja añadida de que el bot responderá ininterrumpidamente a cualquier hora del día.

\newpage

\section{Estructura de la memoria}

La memoria utiliza la siguiente estructura:

\begin{itemize}
	\tightlist
	\item
	\textbf{Introducción:} resumen del contenido del trabajo: problema, proceso y solución. Estructura de la memoria y listado de los materiales adjuntos.
	\item
	\textbf{Objetivos del proyecto:} listado de objetivos generales y técnicos del proyecto.
	\item
	\textbf{Conceptos teóricos:} explicación de los conceptos teóricos necesarios para comprender el proyecto.
	\item
	\textbf{Técnicas y herramientas:} listado de técnicas y herramientas analizadas para el desarrollo. Comparativa y justificación de la elección.
	\item
	\textbf{Aspectos relevantes del desarrollo del proyecto:} resumen de las distintas fases del desarrollo. Análisis de resultados obtenidos por el programa desarrollado.
	\item
	\textbf{Trabajos relacionados:} estado del arte en el campo de los \textit{chatbots} en la UBU y otras unióndades españolas.
	\item
	\textbf{Conclusiones y Líneas de trabajo futuras:} conclusiones extraídas de la elaboración del proyecto y pasos a seguir en la futura explotación y mejora.
\end{itemize}


La estructura de los anexos es la siguiente:

\begin{itemize}
	\tightlist
	\item
	\textbf{Plan de Proyecto Software:} planificación temporal y estudio de viabilidad económica y legal del proyecto.
	\item
	\textbf{Especificación de requisitos:} fase de análisis: objetivos generales, catálogo de requisitos y especificación de requisitos.
	\item
	\textbf{Especificación de diseño:} diseño de datos, diseño procedimental y diseño arquitectónico.
	\item
	\textbf{Documentación técnica de programación:} estructura de directorios, manual del programador, instalación del proyecto y pruebas del sistema.
	\item
	\textbf{Documentación de usuario:} manual para que el usuario sea capaz de utilizar de manera adecuada el \textit{chatbot} en sus distintas integraciones
\end{itemize}


\section{Materiales adjuntos}

Los materiales adjuntos al proyecto son los siguientes:

\begin{itemize}
	\tightlist
	\item
	Proyecto de Dialogflow para el \text{Chatbot} de la modalidad online.
	\item
	Proyecto de Dialogflow para el \text{Chatbot} para la modalidad presencial.
	\item
	Listado de preguntas y respuestas para la modalidad online.
	\item
	Listado de preguntas y respuestas para la modalidad presencial.
	\item
	Código para la integración HTML en UBUVirtual modalidad online.
	\item
	Código para la integración HTML en UBUVirtual modalidad presencial.

\end{itemize}

Están disponibles a través de Internet en el siguiente recurso:
\begin{itemize}
	\tightlist
	\item
	Repositorio del proyecto en Github \cite{repositorioGithub}.
\end{itemize}





\capitulo{2}{Objetivos del proyecto}

Estos son los distintos objetivos con los que se ha realizado el proyecto.

\section{Objetivos generales}\label{objetivos-generales}

\begin{itemize}
	\tightlist
	\item
	Estudiar y evaluar las distintas plataformas y \textit{frameworks} que permiten trabajar con \emph{chatbots}, seleccionando la más adecuada de entre aquellas gratuitas y con soporte para castellano.
	\item
	Desarrollar un \emph{chatbot} que permita interaccionar con los estudiantes y darle respuestas a sus dudas con respecto a la asignatura del Trabajo de Fin de Grado para su versión online.
	\item
	Desarrollar un segundo \emph{chatbot} para la versión presencial con el objetivo de recibir mayor número de logs.
	\item
	Integrarlos en la plataforma de UBUVirtual, bajo las restricciones dadas por la Universidad en cuanto a la edición del código fuente de la web.
	\item 
	Integrar también uno de ellos en otra plataforma.
	\item
	Implementar un sistema de \emph{logs} que permita monitorizar las interacciones con los usuarios.
	\item
	Explotar los datos obtenidos en los \emph{logs}, para aplicarles un análisis numérico del que extraer información. 
\end{itemize}


\section{Objetivos técnicos}\label{objetivos-tecnicos}

\begin{itemize}
	\tightlist 
	\item
	Implementar los \emph{chatbots} en la plataforma Dialogflow.
	\item
	Utilizar la herramienta Dialogflow Messenger para disponer de las respuestas avanzadas.
	\item 
	Utilizar las \emph{payload response} para dotar a nuestro chatbot de funcionalidades más avanzadas.
	\item
	Obtener el código HTML para su integración.
	\item
	Generar el código CSS que le de una apariencia estéticamente agradable al integrarlo en UBUVirtual.
	\item
	Generar un Moodle propio para testear la integración.
	\item
	Integrar los \emph{chatbots} en UBUVirtual.
	\item
	Implementar e integrar el \emph{chatbot} de modalidad online en Slack.
	\item
	Monitorizar las conversaciones de los usuarios.
	\item
	Realizar un estudio numérico de los datos obtenidos en los \textit{logs}.
	\item
	Utilizar el control de versiones mediante GitHub.
	
	
	
\end{itemize}
\capitulo{3}{Conceptos teóricos}

En aquellos proyectos que necesiten para su comprensión y desarrollo de unos conceptos teóricos de una determinada materia o de un determinado dominio de conocimiento, debe existir un apartado que sintetice dichos conceptos.

Algunos conceptos teóricos de \LaTeX \footnote{Créditos a los proyectos de Álvaro López Cantero: Configurador de Presupuestos y Roberto Izquierdo Amo: PLQuiz}.

\section{Secciones}

Las secciones se incluyen con el comando section.

\subsection{Subsecciones}

Además de secciones tenemos subsecciones.

\subsubsection{Subsubsecciones}

Y subsecciones. 


\section{Referencias}

Las referencias se incluyen en el texto usando cite \cite{wiki:latex}. Para citar webs, artículos o libros \cite{koza92}.


\section{Imágenes}

Se pueden incluir imágenes con los comandos standard de \LaTeX, pero esta plantilla dispone de comandos propios como por ejemplo el siguiente:

\imagen{escudoInfor}{Autómata para una expresión vacía}



\section{Listas de items}

Existen tres posibilidades:

\begin{itemize}
	\item primer item.
	\item segundo item.
\end{itemize}

\begin{enumerate}
	\item primer item.
	\item segundo item.
\end{enumerate}

\begin{description}
	\item[Primer item] más información sobre el primer item.
	\item[Segundo item] más información sobre el segundo item.
\end{description}
	
\begin{itemize}
\item 
\end{itemize}

\section{Tablas}

Igualmente se pueden usar los comandos específicos de \LaTeX o bien usar alguno de los comandos de la plantilla.

\tablaSmall{Herramientas y tecnologías utilizadas en cada parte del proyecto}{l c c c c}{herramientasportipodeuso}
{ \multicolumn{1}{l}{Herramientas} & App AngularJS & API REST & BD & Memoria \\}{ 
HTML5 & X & & &\\
CSS3 & X & & &\\
BOOTSTRAP & X & & &\\
JavaScript & X & & &\\
AngularJS & X & & &\\
Bower & X & & &\\
PHP & & X & &\\
Karma + Jasmine & X & & &\\
Slim framework & & X & &\\
Idiorm & & X & &\\
Composer & & X & &\\
JSON & X & X & &\\
PhpStorm & X & X & &\\
MySQL & & & X &\\
PhpMyAdmin & & & X &\\
Git + BitBucket & X & X & X & X\\
Mik\TeX{} & & & & X\\
\TeX{}Maker & & & & X\\
Astah & & & & X\\
Balsamiq Mockups & X & & &\\
VersionOne & X & X & X & X\\
} 

\capitulo{4}{Técnicas y herramientas}

\section{Plataformas para la construcción de Chatbots}

En esta sección vamos a analizar algunas de las herramientas más extendidas para el desarrollo de chatbots. Dialogflow queda exento de esta sección ya que al ser la plataforma elegida será desarrollada más ampliamente en el siguiente apartado. \\

\imagenMediana{logosPlataformas}{Logos plataformas analizadas.}

\newpage
\subsection{Rasa \cite{RasaWeb}} 

Rasa es un framework de aprendizaje automático y código abierto para la creación de chatbots. Está escrito en Python y sirve para desarrollar los chatbots basados tanto en texto como voz. \cite{rasa}

Está compuesto por:
\begin{itemize}
	\item 
	NLU - Comprensión del lenguaje natural. Es el componente que se encarga de procesar la entrada de texto o voz recibida y convertirla en un valor que la máquina pueda entender. 
	\item 
	Core - Es el componente que en base a unos algoritmos y una entrada recibida, decide la acción que el agente debe tomar para continuar la conversación de manera más adecuada. Para ello asigna un valor a cada posible respuesta que se corresponde a la similitud encontrada entre el mensaje que dispara esa respuesta y la entrada recibida. 
\end{itemize}

Por medio de librerías, se puede adaptar el NLU a cualquier idioma.

Las desventajas de esta herramienta son su curva de aprendizaje, la carencia de \textit{hosting} en la nube, no dispone de integraciones \textit{out of box} (que no requieren de instalación) y no dispone en su versión de gratuita de interfaz gráfica para la construcción del chatbot.
\newpage

\subsection{Amazon Lex \cite{AmazonLexWeb}}

Amazon Lex es un servicio de Amazon para la creación de chatbots, basados tanto en texto como voz.
Utiliza el mismo motor que Amazon Alexa, un potente asistente virtual lanzado por Amazon en 2013, y actualmente el más popular del mundo con un 62\% de la cuota de mercado en 2017.\cite{alexa}

Ofrece NLU para la comprensión del lenguaje natural y ASR para el reconocimiento de voz. \\

Sus principales características son:
\begin{itemize}
	\item Dispone de una interfaz gráfica web para el desarrollo, lo que reduce la curva de aprendizaje.
	\item Permite integraciones con servicios de la plataforma AWS.
	\item Soporte para diez idiomas, castellano incluido.
	\item Integraciones \textit{out of box} con Slack y Facebook, entre otras.
\end{itemize}  

Como principal desventaja, a partir del segundo año pasa a ser de pago, por lo que no cumple la principal restricción del proyecto. 
\newpage

\subsection{Watson Assistant \cite{WatsonAssistant}} 

Watson Assistant es un producto de IBM para el desarrollo de chatbots conversacionales.
Permite el salto de un nodo de la conversación a otro, lo cual en nuestra herramienta elegida (Dialogflow) no es posible, teniendo que duplicar diálogos en ocasiones.

Sus principales características son:
\begin{itemize}
	\item Interfaz gráfica web para el desarrollo, sencilla e intuitiva.
	\item Soporte para trece idiomas, castellano incluido.
	\item Genera un \textit{script} para una sencilla integración web.
	\item Fácil integración con Slack, Whatsapp y Facebook, entre otras.
\end{itemize}

Dispone de una versión gratuita, la cual podría cubrir todas las necesidades de nuestro proyecto. Esta versión limita a mil usuarios al mes, lo cual no supone un problema para el espectro de este proyecto; ante una posible futura expansión del asistente a un mayor número de asignaturas o incluso la página oficial de la UBU, podría llegar a suponer un problema. \\

Siendo esta una herramienta que cumple todos los requisitos de nuestro proyecto, se ha realizado una comparativa con Dialogflow para determinar qué plataforma era más recomendable utilizar.
A favor de Watson se ha detectado la ventaja del salto de nodos comentada anteriormente, pero en su contra, todas estas desventajas:
\begin{itemize}
	\item Menos idiomas e integraciones que Dialogflow.
	\item No permite diálogos multidioma, lo cual en una futura actualización inclusiva para los estudiantes Erasmus podría ser beneficioso; ya que junto al punto anterior se podría dar respuesta en mayor número de idiomas.
	\item Está enfocado para desarrollos más complejos que el de nuestro Q/A, para lo cual Dialogflow está mejor preparado.
\end{itemize}
 
 
\newpage

\subsection{Chatcompose \cite{ChatCompose}} 

Chatcompose es una plataforma web para el desarrollo de chatbots.
Es una herramienta muy potente por su rápida personalización, implementación y despliegue en múltiples plataformas.  
Permite añadir \textit{rich responses} de manera automática, sin necesidad de ningún tipo de código o lenguaje especial, simplemente con las herramientas que nos proporciona su interfaz gráfica.

A diferencia del resto de herramientas, proporciona \textit{live chat}, en el que un agente humano puede tomar parte de la conversación. Esta funcionalidad no resulta interesante para nuestra aplicación, ya que no se pretende que ningún miembro de la UBU responda las preguntas, sino que se delega todo en el chatbot.

Sus principales ventajas son:
\begin{itemize}
	\item Soporte para 32 idiomas, castellano incluido.
	\item Única herramienta estudiada con soporte propio para catalán.
	\item Integración con Slack, Telegram, Facebook, Whatsapp, SMS y Line entre otras.
	\item Sencilla y ampliamente personalizable. 
	
\end{itemize}

Es una herramienta muy interesante, pero en su contra tiene que la versión gratuita tiene muchísimas limitaciones, la principal es que únicamente permite desarrollar un chatbot.
Para poder hacer un uso efectivo de la plataforma deberíamos recurrir a la versión de pago.\\

Además, Dialogflow está mejor preparado para conversaciones abiertas, por su fácil y rápida implementación de preguntas y respuestas. Chatcompose sería más eficaz si se guiase al usuario a través de un flujo de conversación, ya que permite implementar las \textit{chip response} sin necesidad de enviar la información en formato JSON como requiere Dialogflow. 
\newpage

\subsection{Otras plataformas}

Las plataformas analizadas representan un pequeño número de la inmensa cantidad de plataformas para el desarrollo de chatbots que existen en la actualidad. Se han seleccionado esas plataformas debido a que en la fecha de inicio del proyecto contaban con soporte nativo para castellano y eran algunas de las soluciones más populares. No obstante, cabe aclarar que existen diversas alternativas que también cumplen dicho criterio y han quedado fuera porque el objetivo de este trabajo no es el de analizar todas las posibles soluciones -que serían muchas-, sino realizar un estudio más exhaustivo de cinco de ellas. 

A continuación se resumen otras alternativas:

\begin{itemize}
	\item \textbf{Microsoft Bot Framework}: también llamado Azure Bot Service, es una solución de Microsoft. Su principal fortaleza es la integración en el ecosistema tecnológico de Microsoft, incluyendo Office y Teams como principales integraciones de las que carecen otras plataformas. Cuenta con soporte para castellano, pero la versión gratuita está limitada a 10000 mensajes mensuales. \cite{MicrosoftBotFramework}
	\item \textbf{Chatfuel}: Es una herramienta muy popular, pero únicamente dispone de integración con Facebook Messenger e Instagram. \cite{ChatFuel}
	\item \textbf{Wit.Ai}: Cuenta con soporte para castellano y es gratuita, aunque con limitaciones. Su punto flojo es la poca documentación e información de la que se dispone. \cite{Wit.Ai}
	\item \textbf{Flow XO}: Ofrece más de 100 integraciones, incluyendo YouTube, GitHub, LinkedIn, Office y MySQL entre otras. No dispone de versión gratuita. \cite{FlowXO}
	
\end{itemize}


\newpage

\subsection{Tabla comparativa}

A continuación se muestra la comparativa entre las distintas plataformas estudiadas en los aspectos que resultan más relevantes para el proyecto.


\tablaSmall{Plataformas de desarrollo de Chatbots.}{l l l l l l}{comparativaplataformas}
{ Plataforma & Castellano & UI Based & Slack Integration & Hosting & Gratuita \\}{ 
	Dialogflow & X & X & X & X & X\\
	Rasa & X &  &  &  & X\\
	Amazon Lex & X & X & X & X & \\
	IBM Watson & X & X & X & X & X\\
	Chatcompose & X & X & X & X & \\
} 

\capitulo{5}{Aspectos relevantes del desarrollo del proyecto}

Este apartado pretende recoger los aspectos más interesantes del desarrollo del proyecto, comentados por los autores del mismo.
Debe incluir desde la exposición del ciclo de vida utilizado, hasta los detalles de mayor relevancia de las fases de análisis, diseño e implementación.
Se busca que no sea una mera operación de copiar y pegar diagramas y extractos del código fuente, sino que realmente se justifiquen los caminos de solución que se han tomado, especialmente aquellos que no sean triviales.
Puede ser el lugar más adecuado para documentar los aspectos más interesantes del diseño y de la implementación, con un mayor hincapié en aspectos tales como el tipo de arquitectura elegido, los índices de las tablas de la base de datos, normalización y desnormalización, distribución en ficheros3, reglas de negocio dentro de las bases de datos (EDVHV GH GDWRV DFWLYDV), aspectos de desarrollo relacionados con el WWW...
Este apartado, debe convertirse en el resumen de la experiencia práctica del proyecto, y por sí mismo justifica que la memoria se convierta en un documento útil, fuente de referencia para los autores, los tutores y futuros alumnos.

\capitulo{6}{Trabajos relacionados}

El primer bot conversacional fue creado entre 1964 y 1966 por Joseph Weizenbaum en el Laboratorio de Ciencias de la Computación e Inteligencia Artificial del Instituto de Tecnología de Massachusetts \cite{wiki:eliza}. Simulaba ser un psicólogo y se creó con el propósito de demostrar la superficialidad de las conversaciones entre humanos y máquinas.

Desde entonces han evolucionado, especialmente en los últimos años y aún más tras la pandemia de Covid-19, reflejo de ello es un estudio de \textit{Markets And Markets} \cite{ChatbotGrowth} que apunta un incremento previsto en el volumen de negocio de los chatbots de 4.8 billones de dólares en 2020 a 13.9 billones en 2025.

Las finalidades de los chatbots son muy variadas: responder preguntas frecuentes, ofrecer soporte al cliente, encuestas, reservas, etc. En este apartado vamos a analizar algunos trabajos relacionados. Al tratarse de un campo tan amplio vamos a focalizar la búsqueda con una serie de criterios. El primer criterio es que únicamente vamos a tener en cuenta los chatbots programados en castellano, ya que el idioma es un aspecto relevante al ser en torno a lo que giran los NLP.
El segundo criterio va a ser el del fin; solo vamos a buscar chatbots implementados en universidades españolas con el fin de responder dudas frecuentes de estudiantes, siendo este el mismo propósito que el de nuestro proyecto.
\\No obstante, en primer lugar vamos a analizar los proyectos de la Universidad de Burgos que tengan alguna relación con los chatbots, para tener una visión de la situación actual de la UBU en esta tecnología.

\newpage


\section{Proyectos de la Universidad de Burgos}

\subsection{\emph{UBUAssistant}}

\imagenPequena{UBUAssistant}{UBUAssistant logo}

Proyecto de Daniel Santidrián Alonso continuado por Carlos González Calatrava en 2018 tutorizado por Pedro Renedo Fernández.
Se trata de un asistente virtual para Android que cuenta con asistente de voz. El objetivo es facilitar las búsquedas en la web de la UBU al realizarlas desde un smartphone.
La aplicación está alojada en un servidor de Azure que ofrece una suscripción gratuita para estudiantes, aunque con ciertas limitaciones. Utiliza el algoritmo de Razonamiento Basado en Casos. \cite{UBUAssistant}


\subsection{\emph{Chatbot for Tourist Recommendations}}

\imagenPequena{touristBot}{Tourist Bot logo}
Proyecto de Jasmin Wellnitz en 2017 tutorizado por el Dr. Bruno Baruque Zanon. 
Chatbot implementado en Telegram actualmente fuera de servicio cuya finalidad era dar recomendaciones turísticas en inglés.
Utilizó la misma tecnología que este proyecto, en aquel entonces se llamaba API.AI y era el mismo NLP que es ahora Dialogflow tras ser comprado en 2016 por Google que optó al año siguiente por cambiarle el nombre.
El proyecto esta alojado en la \textit{Platform as a Service} Heroku. Utiliza la versión de prueba gratuita pero tiene limitaciones en cuanto al uso mensual. \cite{ChatbotTourist}


\subsection{\emph{UBUVoiceAssistant}}

\imagenPequena{UBUVoiceAssistant}{UBUVoiceAssistant logo}

Proyecto de Álvaro Delgado Pascual tutorizado por el Dr. Raúl Marticorena Sanchez.
Se trata de una aplicación que por medio de un asistente de voz permite al usuario obtener información sobre una plataforma de Moodle. Está desarrollado en Python y utiliza el asistente de voz Mycroft. \cite{UBUVoiceAssistant}

\newpage
\section{Proyectos externos}

\subsection{\emph{Chatbot Lola}}

\textit{Lola} es un chatbot de la Universidad de Murcia creado en 2018 cuya finalidad es la de resolver dudas a nuevos estudiantes acerca del proceso de admisión y matrícula.\\
Lola es un proyecto del Servicio de  Información Universitario (SIU) que se ha desarrollado con el Área de Tecnologías de la Información y de las Comunicaciones Aplicadas de la Universidad de Murcia (ATICA) y la empresa 1 million bot. \cite{Lola}\\

\imagenMediana{Lola}{Logo y funcionalidades de Lola}

Este mismo chatbot está implementado por la misma empresa en otras universidades: Universidad de Zaragoza, Complutense de Madrid y Universidad Politécnica de Valencia entre otras. Todos tienen el mismo funcionamiento aunque utilizan nombres comerciales distintos, siendo siempre un nombre de mujer.

\newpage
\section{Tabla comparativa del proyecto} 

En la tabla \ref{tabla:comparativaproyectos} se muestra una comparativa de las características necesarias para este proyecto que aportan cada uno de los proyectos analizados. 

\tablaSmall{Plataformas de desarrollo de Chatbots.}{l l l l l l}{comparativaproyectos}
{ Proyecto & UBUChatbot & UBUAssistant & Tourist Chatbot & UBUVoice & Lola \\}{ 
	Gratuito & X & X & X & X & X\\
	Chatbot QA & X &  &  & & X\\
	Integración Web & X & X &  & X & X\\
	Integración Slack & X &  &  & & \\ 
	Castellano & X & X &  & X & X\\
} 
\capitulo{7}{Conclusiones y Líneas de trabajo futuras}

En este apartado se reflejan las conclusiones sacadas de la realización del proyecto así como una serie de líneas de trabajo extraídas a partir de ellas y que pueden servir como directrices para futuras mejoras y continuidad de la explotación del chatbot.

\section{Conclusiones}
\begin{itemize}
	\tightlist
	\item
	El objetivo del proyecto se ha cumplido, el chatbot resuelve con éxito más del 50\% de las preguntas planteadas por los alumnos, reduciendo con esto de una manera notable el número de consultas realizadas por otros medios al personal docente de la asignatura.
	Además, no se ha incurrido en ningún coste ni introducido ninguna cuenta de facturación en ninguna de las herramientas ni plataformas utilizadas.
	\newline
	\item
	La integración con una plataforma distinta al portal de UBUVirtual se ha realizado correctamente. La plataforma elegida ha sido Slack, aunque también se podía haber realizado con Whatsapp (aunque no de manera gratuita), Facebook Messenger o Telegram, entre otras. 
	
	Esta integración por si misma no aporta valor al proyecto por el momento, pero en futuras versiones del chatbot que implementen nuevas funcionalidades sí podría ser de utilidad.
	\newline
	\item
	Dialogflow ha resultado ser una solución óptima para este proyecto: un NLP potente, gratuita y con una curva de aprendizaje muy pequeña que facilitará la continuidad del chatbot en manos de futuros estudiantes.
	\newline
	\item
	El éxito de la conversación depende en gran medida del correcto uso por parte del estudiante, si sabe formular las preguntas de manera precisa y concisa va a obtener muchos mejores resultados que alguien que formula largas preguntas con un lenguaje similar al que utilizaría si estuviese hablando con un humano. Por esta parte, al tratarse los usuarios de estudiantes de Ingeniería Informática con ciertos conocimientos sobre la limitación de la Inteligencia Artificial, en la mayoría de casos no se da este problema. 
	 
	Además, con la adición de indicaciones de cómo formular las preguntas por parte de nuestro chatbot en caso de mensajes no reconocidos, se ha mejorado en este aspecto.
	\newline
	\item
	Los chatbots son una herramienta en auge, los NLP son cada vez más potentes y sin duda es una tecnología que ha venido para quedarse. Cada vez son más las webs que cuentan con un chatbot y la Universidad de Burgos no puede quedarse atrás en este aspecto. 
	
\end{itemize}

\newpage
\section{Líneas de trabajo futuras}
\begin{itemize}
	\tightlist
	\item
	Integración del chatbot en la página principal de la asignatura. Con esto se obtendrá una mayor visibilidad por parte de los potenciales usuarios, lo que elevará el número de interacciones, reduciendo el número de preguntas que realicen a los profesores. Además, esto nos proporcionará mayor número de \textit{logs}. 
	\newline
	\item
	Actualización y mejora de nuestro NLP con nuevas preguntas y frases de entrenamiento obtenidas a partir de los nuevos \textit{logs} recibidos. El objetivo debe ser continuar mejorando la tasa de éxito de respuesta. 
	\newline
	\item
	Añadir una encuesta al finalizar la conversación en la que se reciba feedback por parte del usuario. Se implementará en forma de preguntas sobre la satisfacción y sugerencias de mejora.\\
	Con esto se obtendrá información más objetiva sobre el grado de satisfacción y mejoras deseadas por parte de nuestros usuarios.
	\newline
	\item
	Actualización a proyecto con cuenta de pago asociada que permita la implementación de código para nuestros \textit{webhooks}. Implementar por código un control de \textit{fallbacks} que permita redirigir la conversación a otro medio una vez se produzcan errores como tres mensajes seguidos sin contestación o la entrada reiterada a un mismo \textit{intent}.
	\newline
	\item 
	Estudiar el posible interés en realizar consultas por código a fuentes externas.
	\newline
	\item
	Analizador léxico para extraer información que permita reconocer patrones de comportamiento.
	\newline
	\item 
	Agente multilingüe para responder preguntas a estudiantes Erasmus o de otros programas de intercambio.
	\newline
	\item
	Implementación de chatbot para la página web principal de la Universidad de Burgos y otras asignaturas.
	\newline
	\item
	Estudiar los posibles beneficios de continuar la implementación en Slack y/o otras plataformas distintas a UBUVirtual.
	
\end{itemize}


\bibliographystyle{ieeetr}
\bibliography{bibliografia}

\end{document}
